\begin{ex}
Số 75 842 có chữ số hàng nghìn là bao nhiêu?
\shortans{5}
\loigiai{
Chữ số hàng nghìn là 5.
}
\end{ex}

\begin{ex}
Số 84 026 là số chẵn hay lẻ? (Chẵn ghi 0, lẻ ghi 1).
\shortans{0}
\loigiai{
Tận cùng là 6 nên là số chẵn (0).
}
\end{ex}

\begin{ex}
Số liền sau của 69 999 là số nào?
\shortans{70000}
\loigiai{
$69999 + 1 = 70000$.
}
\end{ex}

\begin{ex}
Số liền trước của 100 000 là số nào?
\shortans{99999}
\loigiai{
$100000 - 1 = 99999$.
}
\end{ex}

\begin{ex}
Tính giá trị biểu thức $a + b \times 2$ với $a = 15, b = 20$.
\shortans{55}
\loigiai{
$15 + 40 = 55$.
}
\end{ex}

\begin{ex}
Tính giá trị biểu thức $(x - y) : 3$ với $x = 80, y = 20$.
\shortans{20}
\loigiai{
$60 : 3 = 20$.
}
\end{ex}

\begin{ex}
Tính: $45000 + 35000 = \dots$
\shortans{80000}
\loigiai{
$45000 + 35000 = 80000$.
}
\end{ex}

\begin{ex}
Tính: $90000 - 45000 = \dots$
\shortans{45000}
\loigiai{
$90000 - 45000 = 45000$.
}
\end{ex}

\begin{ex}
Tính: $15000 \times 4 = \dots$
\shortans{60000}
\loigiai{
$15000 \times 4 = 60000$.
}
\end{ex}

\begin{ex}
Tính: $80000 : 4 = \dots$
\shortans{20000}
\loigiai{
$80000 : 4 = 20000$.
}
\end{ex}

\begin{ex}
Làm tròn số 64 700 đến hàng nghìn được số nào?
\shortans{65000}
\loigiai{
Làm tròn lên 65 000.
}
\end{ex}

\begin{ex}
Làm tròn số 82 400 đến hàng chục nghìn được số nào?
\shortans{80000}
\loigiai{
Làm tròn thành 80 000.
}
\end{ex}

\begin{ex}
Trong các số $456, 789, 120, 345$, có bao nhiêu số lẻ?
\shortans{2}
\loigiai{
789 và 345 (có 2 số).
}
\end{ex}

\begin{ex}
Số lẻ nhỏ nhất có 5 chữ số là số nào?
\shortans{10001}
\loigiai{
Số đó là 10001.
}
\end{ex}

\begin{ex}
Số chẵn lớn nhất có 5 chữ số là số nào?
\shortans{99998}
\loigiai{
Số đó là 99998.
}
\end{ex}

\begin{ex}
Một cửa hàng có 3 thùng dầu, thùng 1 có 20 lít, thùng 2 có gấp đôi thùng 1, thùng 3 có 30 lít. Cả 3 thùng có bao nhiêu lít dầu?
\shortans{90}
\loigiai{
$20 + 40 + 30 = 90$ lít.
}
\end{ex}

\begin{ex}
Tìm $x$: $x + 18000 = 50000$.
\shortans{32000}
\loigiai{
$x = 50000 - 18000 = 32000$.
}
\end{ex}

\begin{ex}
Tìm $x$: $x : 3 = 12000$.
\shortans{36000}
\loigiai{
$x = 12000 \times 3 = 36000$.
}
\end{ex}

\begin{ex}
Tính giá trị biểu thức: $100000 - 25000 \times 3$.
\shortans{25000}
\loigiai{
$100000 - 75000 = 25000$.
}
\end{ex}

\begin{ex}
Tính nhanh: $25000 + 47000 + 25000$.
\shortans{97000}
\loigiai{
$50000 + 47000 = 97000$.
}
\end{ex}

\begin{ex}
Hiệu của số chẵn lớn nhất có 2 chữ số và số chẵn bé nhất có 2 chữ số là bao nhiêu?
\shortans{88}
\loigiai{
$98 - 10 = 88$.
}
\end{ex}

\begin{ex}
Hiệu của số lẻ lớn nhất có 2 chữ số và số lẻ bé nhất có 2 chữ số là bao nhiêu?
\shortans{88}
\loigiai{
$99 - 11 = 88$.
}
\end{ex}

\begin{ex}
Cho chu vi hình chữ nhật là $P = (a + b) \times 2$. Với $a = 25$ m, $b = 15$ m, tính $P$ theo mét.
\shortans{80}
\loigiai{
$P = (25 + 15) \times 2 = 80$ m.
}
\end{ex}

\begin{ex}
Một hình vuông có cạnh $a = 12$ cm. Chu vi hình vuông là bao nhiêu xăng-ti-mét?
\shortans{48}
\loigiai{
$12 \times 4 = 48$ cm.
}
\end{ex}

\begin{ex}
Tổng của 3 số tự nhiên liên tiếp là 30. Số ở giữa là bao nhiêu?
\shortans{10}
\loigiai{
Số ở giữa là $30 : 3 = 10$.
}
\end{ex}
